\part{Anpassen}

\chapter{Eine Zahl statt zwanzig}

\section{Warum Schichten nicht reichen}

\begin{merksatz}
\boxtitel{Die Anpassungsformel in der Praxis}
Satz 6.2 verlangt, über alle Werte von $Z$ zu summieren. Bei zwei Stufen ist das leicht.

\begin{center}
{\small
\renewcommand{\arraystretch}{1.48}
\begin{tabular}{@{}P{3.0cm}P{2.6cm}P{2.6cm}@{}}
\toprule
\textbf{Störgrößen} & \textbf{Schichten} & \textbf{je Schicht}\\
\midrule
$1$ binär & $2$ & $600$\\
$5$ binär & $32$ & $38$\\
$10$ binär & $1024$ & $1$\\
$3$ stetig & $\infty$ & $0$\\
\bottomrule
\end{tabular}}
\end{center}
\kursiv{Bei $1200$ Patienten.}

\begin{center}
\begin{minipage}{0.9\textwidth}
\textbf{Ab etwa fünf Störgrößen ist keine Schicht mehr besetzt --- der Fluch der Dimension aus Band 22, Kapitel~10.}
\end{minipage}
\end{center}
\end{merksatz}

\begin{definition}
\boxtitelm{Definition 7.1 --- Behandlungsneigung}
Die \Begriff{Behandlungsneigung} eines Patienten ist die Wahrscheinlichkeit, mit der er die Behandlung erhält:
\[ e(\vek z) = \Pr[X=1 \mid Z=\vek z] . \]
\end{definition}

\begin{merksatz}
\boxtitel{Satz 7.2 --- Rosenbaum und Rubin, 1983}
Genügt $Z$ dem Hintertürkriterium, so genügt auch die einzelne Zahl $e(Z)$:
\[ \bigl(Y(1),Y(0)\bigr) \perp X \ \mid\ Z \quad\Longrightarrow\quad \bigl(Y(1),Y(0)\bigr) \perp X \ \mid\ e(Z) . \]

\schritt{Beweis}
\kursiv{Erstens:} Bei gegebenem $e(Z)$ ist $\Pr[X=1 \mid e(Z)] = e(Z)$, denn
\[ \Pr[X=1 \mid e(Z)] = \mathbb{E}\bigl[\Pr[X=1\mid Z] \ \big|\ e(Z)\bigr] = \mathbb{E}[e(Z)\mid e(Z)] = e(Z) . \]

\kursiv{Zweitens:} Für die potentiellen Ergebnisse gilt ebenso
\[ \Pr\bigl[X=1 \mid Y(1),Y(0),\ e(Z)\bigr] = \mathbb{E}\bigl[\Pr[X=1\mid Y(1),Y(0),Z] \ \big|\ Y(1),Y(0),e(Z)\bigr] . \]
Nach Voraussetzung ist der innere Ausdruck gleich $e(Z)$, unabhängig von den potentiellen Ergebnissen.

\kursiv{Drittens:} Beide Ausdrücke sind gleich $e(Z)$ --- also hängt die Behandlungswahrscheinlichkeit bei gegebenem $e(Z)$ nicht von den potentiellen Ergebnissen ab. $\blacksquare$

\begin{center}
\begin{minipage}{0.9\textwidth}
\textbf{Zwanzig Störgrößen lassen sich auf eine einzige Zahl zusammenziehen --- ohne etwas zu verlieren.}
\end{minipage}
\end{center}

\begin{center}
\begin{minipage}{0.88\textwidth}
\kursiv{Das ist der wichtigste praktische Satz dieses Bandes. \textbf{Er verwandelt ein hochdimensionales Problem in ein eindimensionales}, und die Zahl $e(Z)$ hat eine klare Bedeutung: Zwei Patienten mit gleicher Behandlungsneigung sind, was die Zuteilung angeht, austauschbar --- als hätte man bei ihnen dieselbe Münze geworfen.}
\end{minipage}
\end{center}
\end{merksatz}

\section{Vier Wege, sie zu benutzen}

\begin{merksatz}
\boxtitel{Die Verfahren}
\begin{center}
{\small
\renewcommand{\arraystretch}{1.55}
\begin{tabular}{@{}P{2.4cm}P{6.2cm}@{}}
\toprule
\textbf{Schichten} & Fünf Gruppen gleicher Größe nach $e$; in jeder vergleichen, dann mitteln.\\
\textbf{Paaren} & Zu jedem Behandelten den Unbehandelten mit nächstliegendem $e$ suchen.\\
\textbf{Gewichten} & Jeden Patienten mit $1/e$ bzw.\ $1/(1-e)$ gewichten --- Kapitel~8.\\
\textbf{Als Kovariable} & $e$ ins Regressionsmodell aufnehmen. \textcolor{akzent}{Am schwächsten}, weil es eine Modellform unterstellt.\\
\bottomrule
\end{tabular}}
\end{center}
\end{merksatz}

\begin{merksatz}
\boxtitel{Der Leitfall gerechnet}
$20\,000$ Patienten, zwei stetige Störgrößen, wahres Modell bekannt.

\begin{center}
{\small
\renewcommand{\arraystretch}{1.46}
\begin{tabular}{@{}P{1.8cm}P{1.6cm}P{1.6cm}P{1.6cm}P{2.0cm}@{}}
\toprule
\textbf{Schicht} & $\bar e$ & $n_{\text{beh}}$ & $n_{\text{unbeh}}$ & \textbf{Differenz}\\
\midrule
$1$ & $0{,}12$ & $499$ & $3501$ & $-0{,}039$\\
$2$ & $0{,}27$ & $1112$ & $2888$ & $-0{,}081$\\
$3$ & $0{,}42$ & $1732$ & $2268$ & $-0{,}122$\\
$4$ & $0{,}59$ & $2385$ & $1615$ & $-0{,}138$\\
$5$ & $0{,}80$ & $3156$ & $844$ & $-0{,}124$\\
\bottomrule
\end{tabular}}
\end{center}

\begin{center}
{\small
\renewcommand{\arraystretch}{1.5}
\begin{tabular}{@{}P{3.0cm}P{2.4cm}@{}}
\toprule
\textbf{Verfahren} & \textbf{Schätzung}\\
\midrule
naiv & $+0{,}049$\\
fünf Schichten & $-0{,}101$\\
Gewichten & $\mathbf{-0{,}123}$\\
\midrule
wahrer Effekt & $\mathbf{-0{,}122}$\\
\bottomrule
\end{tabular}}
\end{center}

\begin{center}
\begin{minipage}{0.9\textwidth}
\textbf{Der naive Vergleich hat das falsche Vorzeichen. Fünf Schichten beseitigen den größten Teil der Verzerrung, ein Sechstel bleibt.}
\end{minipage}
\end{center}

\begin{center}
\begin{minipage}{0.88\textwidth}
\kursiv{\textbf{Der Rest ist Restverzerrung innerhalb der Schichten}: In Schicht 1 reicht die Neigung von $0{,}02$ bis $0{,}20$, und innerhalb dieser Spanne sind Behandelte noch immer kränker. Zehn oder zwanzig Schichten verringern den Rest, dünnen aber die Besetzung aus --- derselbe Zielkonflikt wie in Band 22, Kapitel~2. Das Gewichten umgeht ihn, indem es gar nicht erst einteilt.}
\end{minipage}
\end{center}
\end{merksatz}

\section{Der Bruchfall}

\begin{warnung}
\boxtitelr{Die Neigung ist geschätzt}
\begin{center}
{\small
\renewcommand{\arraystretch}{1.52}
\begin{tabular}{@{}P{2.6cm}P{6.0cm}@{}}
\toprule
\textbf{Modellfehler} & $e(Z)$ kennt niemand; man schätzt es, meist durch logistische Regression. Ist das Modell falsch, bleibt Verzerrung.\\
\textbf{Kein Ersatz} & Die Neigung fasst nur zusammen, was in $Z$ steht. Fehlt eine Störgröße, hilft sie nicht.\\
\textbf{Gute Balance täuscht} & Dass Behandelte und Unbehandelte nach dem Paaren ähnlich aussehen, beweist nichts über unbeobachtete Größen.\\
\bottomrule
\end{tabular}}
\end{center}

\begin{center}
\begin{minipage}{0.88\textwidth}
\kursiv{Die dritte Zeile erklärt einen verbreiteten Irrtum. \textbf{Eine Tabelle, die nach dem Paaren gleiche Mittelwerte in allen gemessenen Größen zeigt, sieht aus wie eine randomisierte Studie und ist keine.} Die Randomisierung gleicht auch das Ungemessene aus, die Neigung nur das Gemessene --- Kapitel~3, Satz 3.1.}
\end{minipage}
\end{center}

\schritt{Und ein Fehler beim Aufbau des Modells}
\begin{center}
\begin{minipage}{0.88\textwidth}
\kursiv{In das Neigungsmodell gehören Größen, die den \kursiv{Ausgang} beeinflussen --- nicht solche, die nur die Behandlung beeinflussen. \textbf{Letztere sind Instrumente und verstärken nach Kapitel~5 die Restverzerrung.} Wer das Neigungsmodell nach Güte der Vorhersage von $X$ auswählt, wählt also genau falsch: Ein Modell, das die Behandlung perfekt vorhersagt, macht jede Schätzung unmöglich --- es verletzt die Positivität.}
\end{minipage}
\end{center}
\end{warnung}

\section{Übungsaufgaben}

\begin{aufgabe}
\boxtitel{Aufgabe 7.1 --- Schichten von Hand}
Aus den fünf Schichten der Tabelle: Berechnen Sie den geschichteten Schätzer und den Effekt bei den Behandelten.

\loesung
\schritt{Schritt 1 --- Der Durchschnittseffekt}
Jede Schicht enthält $4000$ von $20\,000$ Patienten, also je $20$ Prozent:
\[ \hat\tau = 0{,}2\,(-0{,}039-0{,}081-0{,}122-0{,}138-0{,}124) = 0{,}2\cdot(-0{,}504) = -0{,}101 \]

\schritt{Schritt 2 --- Der Effekt bei den Behandelten}
Nun gewichtet man mit dem Anteil der \kursiv{Behandelten} je Schicht:
\begin{center}
{\small
\renewcommand{\arraystretch}{1.46}
\begin{tabular}{@{}P{1.6cm}P{1.8cm}P{1.8cm}P{2.2cm}@{}}
\toprule
\textbf{Schicht} & $n_{\text{beh}}$ & \textbf{Anteil} & \textbf{Beitrag}\\
\midrule
$1$ & $499$ & $0{,}056$ & $-0{,}0022$\\
$2$ & $1112$ & $0{,}125$ & $-0{,}0101$\\
$3$ & $1732$ & $0{,}195$ & $-0{,}0238$\\
$4$ & $2385$ & $0{,}268$ & $-0{,}0370$\\
$5$ & $3156$ & $0{,}355$ & $-0{,}0441$\\
\midrule
& $8884$ & $1{,}000$ & $\mathbf{-0{,}117}$\\
\bottomrule
\end{tabular}}
\end{center}

\begin{center}
\begin{minipage}{0.9\textwidth}
\textbf{$-0{,}117$ gegen $-0{,}101$: Die Behandelten profitieren stärker, weil sie in den oberen Schichten überwiegen.}
\end{minipage}
\end{center}

\schritt{Schritt 3 --- Die Lehre}
\begin{center}
\begin{minipage}{0.88\textwidth}
\kursiv{\textbf{Dieselben Schichtdifferenzen, zwei verschiedene Zahlen --- je nachdem, über welche Bevölkerung man mittelt.} Das ist Definition 2.3 in der Rechnung: Der Durchschnittseffekt und der Effekt bei den Behandelten unterscheiden sich, sobald die Wirkung zwischen den Schichten schwankt. Welche Zahl man braucht, hängt von der Frage ab --- Leitlinie oder Nachbetrachtung.}
\end{minipage}
\end{center}
\end{aufgabe}

\begin{aufgabe}
\boxtitel{Aufgabe 7.2 --- Wie fein soll man schichten?}
Bei $1200$ Patienten und einem Behandeltenanteil von $50$ Prozent: Wie viele Schichten sind sinnvoll?

\loesung
\schritt{Schritt 1 --- Die Besetzung}
\begin{center}
{\small
\renewcommand{\arraystretch}{1.46}
\begin{tabular}{@{}P{2.0cm}P{2.2cm}P{2.6cm}@{}}
\toprule
\textbf{Schichten} & \textbf{je Schicht} & \textbf{je Gruppe}\\
\midrule
$5$ & $240$ & $120$\\
$10$ & $120$ & $60$\\
$20$ & $60$ & $30$\\
$50$ & $24$ & $12$\\
\bottomrule
\end{tabular}}
\end{center}

\schritt{Schritt 2 --- Die verbleibende Verzerrung}
Eine Faustregel besagt, dass fünf Schichten etwa $90$ Prozent der Verzerrung einer stetigen Störgröße beseitigen, zehn Schichten etwa $95$.

\kursiv{Im Leitfall:} $-0{,}101$ gegenüber $-0{,}122$ sind $83$ Prozent --- etwas weniger, weil die Neigung stark streut.

\schritt{Schritt 3 --- Die Streuung}
Bei $30$ Patienten je Gruppe und Schicht beträgt die Streuung einer Differenz zweier Anteile nahe $0{,}3$ rund
\[ \sqrt{\frac{2\cdot0{,}21}{30}} = 0{,}118 . \]
Über $20$ Schichten gemittelt: $0{,}118/\sqrt{20} = 0{,}026$.

\begin{center}
\begin{minipage}{0.88\textwidth}
\kursiv{\textbf{Die Streuung des Gesamtschätzers wächst durch feinere Schichtung kaum} --- solange jede Schicht beide Gruppen enthält. Bricht das zusammen, so fällt die Schicht ganz aus, und die Schätzung wird verzerrt \kursiv{und} ungenau. Die Grenze liegt also nicht bei der Streuung, sondern bei der Positivität: Man schichtet so fein, wie es die Besetzung zulässt.}
\end{minipage}
\end{center}
\end{aufgabe}

\begin{aufgabe}
\boxtitel{Aufgabe 7.3 --- Was Balance beweist}
Nach dem Paaren stimmen Alter, Vorerkrankungen und Laborwerte in beiden Gruppen überein. Welche Schlüsse sind zulässig?

\loesung
\schritt{Schritt 1 --- Was gezeigt ist}
\begin{center}
{\small
\renewcommand{\arraystretch}{1.5}
\begin{tabular}{@{}P{2.8cm}P{5.8cm}@{}}
\toprule
\textbf{Gezeigt} & Das Neigungsmodell hat die gemessenen Größen ausgeglichen. Es ist also brauchbar angepasst.\\
\textbf{Nicht gezeigt} & Dass es keine weiteren Störgrößen gibt.\\
\bottomrule
\end{tabular}}
\end{center}

\schritt{Schritt 2 --- Warum das kein Zirkel ist}
\begin{center}
\begin{minipage}{0.88\textwidth}
\kursiv{Die Balance prüft das \kursiv{Modell}, nicht die \kursiv{Annahme}. \textbf{Sie ist eine notwendige, keine hinreichende Bedingung}: Wer sie nicht erreicht, hat sicher ein Problem; wer sie erreicht, hat eines der beiden möglichen Probleme ausgeschlossen.}
\end{minipage}
\end{center}

\schritt{Schritt 3 --- Die einzige echte Prüfung}
\begin{center}
{\small
\renewcommand{\arraystretch}{1.5}
\begin{tabular}{@{}P{2.8cm}P{5.8cm}@{}}
\toprule
\textbf{Negative Kontrolle} & Ein Ausgang, auf den die Behandlung nicht wirken kann. Findet man dort einen Effekt, fehlt eine Störgröße.\\
\textbf{Empfindlichkeit} & Wie stark müsste eine verborgene Größe wirken?\\
\textbf{Vergleich} & Gibt es zu derselben Frage eine randomisierte Studie?\\
\bottomrule
\end{tabular}}
\end{center}

\begin{center}
\begin{minipage}{0.88\textwidth}
\kursiv{\textbf{Der dritte Punkt ist die strengste Prüfung und in großem Maßstab durchgeführt worden}: Man bildete mit Beobachtungsdaten Studien nach, deren randomisiertes Ergebnis bekannt war. Bei sorgfältiger Emulation --- Kapitel~11 --- stimmten die Ergebnisse überwiegend überein; bei nachlässiger wichen sie erheblich ab, und zwar in beide Richtungen.}
\end{minipage}
\end{center}
\end{aufgabe}

\begin{geschichte}
\boxtitelg{Ein Aufsatz von acht Seiten}
\Person{Paul Rosenbaum} und \Person{Donald Rubin} veröffentlichten 1983 in \kursiv{Biometrika} den Aufsatz, der Satz 7.2 enthält.

\begin{center}
\begin{minipage}{0.9\textwidth}
\kursiv{Er gehört zu den meistzitierten Arbeiten der Statistik überhaupt. \textbf{Der Gedanke, eine ganze Menge von Störgrößen durch eine einzige Zahl zu ersetzen, war neu und überraschend} --- und der Beweis passt, wie oben, auf eine halbe Seite.}
\end{minipage}
\end{center}

\schritt{Ein früher Einsatz}
\begin{center}
\begin{minipage}{0.9\textwidth}
\kursiv{Rubin wurde in den 1990er Jahren als Gutachter in einem Verfahren gegen Tabakkonzerne herangezogen und benutzte Neigungsverfahren, um die Krankheitskosten des Rauchens zu schätzen. \textbf{Die Methode wurde dadurch weit über die Statistik hinaus bekannt} --- und zugleich zum Gegenstand der Kritik, weil beide Seiten mit denselben Daten verschiedene Zahlen erzeugten.}
\end{minipage}
\end{center}
\end{geschichte}

\begin{forschung}
\boxtitelp{Neigung heute}
\begin{center}
{\small
\renewcommand{\arraystretch}{1.45}
\begin{tabular}{@{}P{2.8cm}P{5.4cm}@{}}
\toprule
\textbf{Maschinell geschätzt} & Die Neigung mit Verfahren aus Band 22 statt logistischer Regression.\\
\textbf{Balance als Ziel} & Gewichte direkt so bestimmen, dass die Mittelwerte übereinstimmen.\\
\textbf{Hochdimensional} & Hunderte von Störgrößen aus Abrechnungsdaten.\\
\bottomrule
\end{tabular}}
\end{center}

\schritt{Zur ersten Zeile}
\begin{center}
\begin{minipage}{0.88\textwidth}
\kursiv{Ein flexibles Lernverfahren schätzt $e(Z)$ genauer --- und schafft ein neues Problem: \textbf{Seine Fehler sind nicht mehr klein genug, damit die einfachen Schätzer gültig bleiben.} Man braucht dann die doppelt robusten Verfahren des nächsten Kapitels, die genau für diesen Fall gebaut sind.}
\end{minipage}
\end{center}
\end{forschung}

\begin{merksatz}
\boxtitel{Was Sie jetzt können}
\begin{center}
\renewcommand{\arraystretch}{1.5}
\begin{tabular}{@{}P{3.0cm}P{5.6cm}@{}}
\toprule
\textbf{Definition 7.1} & Behandlungsneigung $e(\vek z)$\\
\textbf{Satz 7.2} & eine Zahl ersetzt alle Störgrößen\\
\textbf{Vier Wege} & Schichten, Paaren, Gewichten, Kovariable\\
\textbf{Restverzerrung} & innerhalb der Schichten\\
\textbf{Balance} & prüft das Modell, nicht die Annahme\\
\bottomrule
\end{tabular}
\end{center}

\begin{center}
\begin{minipage}{0.88\textwidth}
\kursiv{Das Schichten verschenkt Information, das Gewichten nicht. Das nächste Kapitel führt es aus --- und zeigt ein Verfahren, das selbst dann noch richtig rechnet, wenn eines von zwei Modellen falsch ist.}
\end{minipage}
\end{center}
\end{merksatz}
